\documentclass[../../main]{subfiles}

\begin{document}
\chapter{Álgebra}
\label{chap:matematica:algebra}

\section{Fundamentos}
\label{sec:matematica:algebra:fundamentos}

\begin{defn}[Par Ordenado]
    \label{def:matematica:algebra:fundamentos:par-ordenado}

    É um par de elementos a qual a ordem de ocorrência de cada elemento é
    significante. Geralmente representado por $(a,b)$. Seja $a$ e $b$ elementos
    quaisquer.
\end{defn}

\begin{defn}[Operação Matemática]
    \label{def:matematica:algebra:fundamentos:operacao-matematica}

    É qualquer tipo de procedimento que envolve uma quantidade de elementos de
    um conjunto \cref{def:matematica:teoria-dos-conjuntos:fundamentos:conjunto}
    e resulta em um terceiro elemento de um conjunto.
\end{defn}

\begin{defn}[Operação Binária]
    \label{def:matematica:algebra:fundamentos:operacao-binaria}

    É uma regra que associa dois elementos de um conjunto a um terceiro
    elemento de um conjunto.
\end{defn}

\begin{defn}[Estrutura Algébrica]
    \label{def:matematica:algebra:fundamentos:estrutura-algebrica}

    É um modelo associado a um conjunto de elementos e uma ou mais operações
    binárias que satisfazem um sistema axiomático
    (\cref{def:matematica:logica-matematica:fundamentos:sistema-axiomatico}).
\end{defn}

\begin{defn}[Anel]
    \label{def:matematica:algebra:fundamentos:anel}

    É uma estrutura algébrica que consiste em uma tripla ordenada
    $(A,+,\cdot)$, seja $A$ um conjunto que possui o elemento vazio $\Set{0}$ e
    duas operações binárias definidas em $A$, $+$ e $\cdot$ que satisfazem os
    seguintes axiomas:

    \begin{enumerate}
        \item \textbf{Associatividade de $+$:} $(\forall a,b,c \in A):
            (a+b)+c=a+(b+c)$

        \item \textbf{Comutatividade de $+$:} $(\forall a,b \in A): a+b=b+a$

        \item \textbf{Existência do elemento neutro de $+$:} $(\forall a \in
            A): \Set{I}+a=a+\Set{I}=a$

        \item \textbf{Existência do elemento inverso de $+$:} $(\forall a \in
            A): a+a^{-1}=a^{-1}+a=\Set{I}$
        \item \textbf{Associatividade de $\cdot$:} $(\forall a,b,c \in A):
            (a\cdot b)\cdot c=a\cdot (b\cdot c)$

        \item \textbf{Distributividade de $\cdot$ em relação a $+$ (à esquerda
            e à direita):} $(\forall a,b,c \in A): a\cdot (b+c)=a\cdot b+a\cdot
            c$ e $(\forall a,b,c \in A): a \cdot (b+c)=a \cdot b+a \cdot c
            \wedge (a+b)\cdot c=a\cdot c+b\cdot c$
    \end{enumerate}
\end{defn}

\begin{defn}[Anel Comutativo]
    \label{def:matematica:algebra:fundamentos:anel-comutativo}

    Um anel $(A,+,\cdot)$ \cref{def:matematica:algebra:fundamentos:anel} é
    denominado anel comutativo se além de satisfazer os axiomas de um anel,
    também satisfaz o axioma de comutatividade para a operação $\cdot$, sendo
    $(\forall a,b \in A): a\cdot b=b\cdot a$.
\end{defn}

\begin{defn}[Grupo]
    \label{def:matematica:algebra:fundamentos:grupo}

    Seja $\Set{G}$ um conjunto de elementos e $\*$ uma operação binária em
    $\Set{G}$. O par ordenado $(\Set{G},*)$ é definido um grupo se são
    satisfeitos os seguintes axiomas:

    \begin{enumerate}
        \item \textbf{Fechamento:} Para toda operação binária $a*b=c$, se $a,b
            \in\Set{G}$, então $c \in\Set{G}$.

        \item \textbf{Associatividade:} Para quaisquer elementos $a,b,c
            \in\Set{G}$, $a*(b*c)=(a*b)*c$

        \item \textbf{Existência do elemento neutro:} Existe um elemento
            $\Set{I} \in\Set{G}$ tal que $a*\Set{I}=\Set{I}*a=a$ para qualquer
            elemento $a \in\Set{G}$

        \item \textbf{Existência do elemento inverso:} Dado o elemento neutro
            $\Set{I} \in\Set{G}$, Existe um elemento $a^{-1} \in\Set{G}$ tal
            que $a*a^{-1}=a^{-1}*a=\Set{I}$ para qualquer elemento $a
            \in\Set{G}$
    \end{enumerate}
\end{defn}

\begin{defn}[Corpo]
    \label{def:matematica:algebra:fundamentos:corpo}

    Um anel comutativo $F$
    \cref{def:matematica:algebra:fundamentos:anel-comutativo} é denominado
    corpo se satisfaz a condição: $(\forall a \in
    F\ \backslash\ \{0\})(\exists! b \in F): a\cdot b=\Set{I}$
\end{defn}

\begin{defn}[Espaço Vetorial]
    \label{def:matematica:algebra:fundamentos:espaco-vetorial}

    Se $V$ é dito ser um espaço vetorial sobre o corpo $K$
    \ref{def:matematica:algebra:fundamentos:corpo} se e somente se é munido de
    de duas operações binárias, $+$ e $.$ que satisfazem os seguintes axiomas:

    \begin{enumerate}
        \item \textbf{Associatividade de $+$:} $(\forall a,b,c \in V):
            (a+b)+c=a+(b+c)$

        \item \textbf{Comutatividade de $+$:} $(\forall a,b \in V): a+b=b+a$

        \item \textbf{Existência do elemento neutro de $+$:} $(\forall a \in V,
            \exists! \Set{I} \in V): \Set{I}+a=a+\Set{I}=a$

        \item \textbf{Existência do elemento inverso de $+$:} $(\forall a \in
            V, \exists! a^{-1} \in V): a+a^{-1}=a^{-1}+a=\Set{I}$

        \item \textbf{Associatividade de $.$:} $(\forall a \in V, \alpha, \beta
            \in K): \alpha(\beta.a) = (\alpha.\beta).a $

        \item \textbf{Distributividade de $.$ em relação a $+$:} $(\forall a, b
            \in V, \alpha in K): \alpha.(a+b)=\alpha.a+\alpha.b$

        \item \textbf{Distributividade de $.$ em relação a soma de escalares:}
            $(\forall a \in V, \alpha, \beta \in K):
            (\alpha+\beta).a=(\alpha.a)+(\beta.a)$

        \item \textbf{Existência do elemento neutro de $.$:} $(\forall a \in V,
            \Set{I} \in K): \alpha.a=a$
    \end{enumerate}
\end{defn}

A definição \ref{def:matematica:algebra:fundamentos:espaco-vetorial} foi
baseada em \cite{lima-2012}.
\end{document}
